\chapter{Geometrie, Ansatz und Quadratur}\label{chap:geometry}
\section{Affines Tetraeder und Orientierung}
Mit vier Ecken $\mathbf x_a$ und Referenzkoordinaten
$\boldsymbol\xi=(\xi,\eta,\zeta)$ auf
$\hat K=\{\xi,\eta,\zeta\ge0:\xi+\eta+\zeta\le1\}$ gilt
\begin{equation}\label{eq:affine}
\mathbf x=\mathbf x_0+\mathbf J\boldsymbol\xi,\qquad
\mathbf J=[\mathbf x_1-\mathbf x_0,\mathbf x_2-\mathbf x_0,
\mathbf x_3-\mathbf x_0],\qquad V=|\det\mathbf J|/6.
\end{equation}
Die baryzentrischen Funktionen sind
$L_0=1-\xi-\eta-\zeta$, $L_1=\xi$, $L_2=\eta$, $L_3=\zeta$.
Aus der Kettenregel folgt $\mathbf J^T\nabla_x L_a=\nabla_\xi L_a$.
Der Code löst dieses System; er bildet keine explizite Inverse.
Negative Eingabeorientierung wird mit konsistenter Knotenumordnung korrigiert;
Rohvorzeichen bleiben diagnostizierbar. Bei Tet10 werden alle Kantenknoten
mitpermutiert. Die Feldoperatoren beziehen sich auf die ursprüngliche
Eingabereihenfolge, nicht auf eine unbemerkte Umnummerierung der DOFs.

Das maximale Kantenmaß $\ell_{\max}$ liefert
$q_V=6\sqrt2 V/\ell_{\max}^3$. Die Standardgrenze $q_V\le10^{-12}$ verwirft
entartete oder numerisch unauflösbare Zellen; $q_V=1$ für das reguläre
Tetraeder. Die Größe ist translations-, rotations- und skaleninvariant.
Ein zusätzlicher Konditionsindikator verwendet die zentrierten Eckkoordinaten.
Gute Algebra ersetzt trotzdem keine physikalisch angemessene Netzauflösung.
\index{Tet4}\index{Tet10}\index{Jacobiabbildung}

\section{Lineare und quadratische Interpolation}
Für Tet4 ist $N_a=L_a$; jedes skalare Feld hat vier, jede Verschiebung zwölf
lokale DOFs. Partition der Eins und affine Reproduktion ergeben sich aus
$\sum_a L_a=1$ und $\sum_a L_a\mathbf x_a=\mathbf x$.
Für Tet10 lautet die quadratische nodale Lagrange-Basis
\begin{align}\label{eq:tet10shape}
N_i&=L_i(2L_i-1),&N_{ij}&=4L_iL_j,\\
\nabla N_i&=(4L_i-1)\nabla L_i,&
\nabla N_{ij}&=4(L_i\nabla L_j+L_j\nabla L_i).\nonumber
\end{align}
An den Ecken und Kantenmitten ist $N_a(\mathbf x_b)=\delta_{ab}$; damit folgt
Eindeutigkeit der quadratischen Interpolation. Die Basis ist in
\cite{defelement-p2} dargestellt. Die Toolbox-Reihenfolge ist eigene Konvention:
Ecken 0--3, dann Kanten 01,12,20,03,13,23. Triangle6 verwendet drei Ecken
und 01,12,20. Alle Mittelknoten müssen geometrische Kantenmitten sein
(relative Toleranz $10^{-10}$). Das Feld ist quadratisch, die Geometrie bleibt
affin. Gekrümmte, hybride Tet4/Tet10- oder nichtkonforme Übergänge sind gesperrt.

Mit $\mathbf N=[N_0,\ldots,N_{n-1}]$ gilt
\begin{equation}
T_h=\mathbf N\mathbf T_e,\qquad
\phi_h=\mathbf N\boldsymbol\phi_e,\qquad
\mathbf u_h=\sum_a N_a\mathbf u_a.
\end{equation}
$\mathbf G=[\nabla N_0,\ldots,\nabla N_{n-1}]$ hat Größe 3$\times n$.
Der Engineering-Verzerrungsblock je Knoten ist
\begin{equation}\label{eq:B}
\mathbf B_a=\begin{pmatrix}
N_{a,x}&0&0\\0&N_{a,y}&0\\0&0&N_{a,z}\\
0&N_{a,z}&N_{a,y}\\N_{a,z}&0&N_{a,x}\\N_{a,y}&N_{a,x}&0
\end{pmatrix},\qquad\varepsilon_v=\mathbf B\mathbf u_e.
\end{equation}
$\mathbf B$ hat 6$\times12$ bzw. 6$\times30$ Einträge mit Einheit 1/m.
Tet4-Gradienten sind konstant; Tet10-Gradienten hängen vom Integrationspunkt
ab. Eine einzige Schwerpunktprobe ersetzt daher keine Tet10-Matrixintegration.

\section{Volumen- und Randintegration}
Auf dem Referenztetraeder folgt aus wiederholter Beta-Integration
\begin{equation}\label{eq:simplex-moment}
\int_K\prod_{i=0}^3L_i^{a_i}\,\mathrm dV
=6V\frac{\prod_i a_i!}{(\sum_i a_i+3)!}.
\end{equation}
Die Dreiecksform hat $2A$ und Nenner $(\sum_i a_i+2)!$.
Für $a_i=0$ liefert die Formel $V$ bzw. $A$; für $a_i=1$ an einer Ecke $V/4$.
Sie ist eine unabhängige Quelle für Last- und Massenreferenzen.

Der Tet4-Kern integriert konstante Steifigkeitsintegranden exakt. Seine
Volumen- und Flächenlastregeln vom Grad zwei integrieren affine Lasten exakt,
weil $N$ linear ist. Die Tet10-Referenz nutzt kollabierte Gaußprodukte:
\begin{align}\label{eq:duffy}
(a,b,c)&\mapsto(a,(1-a)b,(1-a)(1-b)c),&&J_D=(1-a)^2(1-b),\\
(a,b)&\mapsto(a,(1-a)b),&&J_D=1-a.\nonumber
\end{align}
Ein eindimensionales $n$-Punkt-Gaußverfahren integriert bis Grad $2n-1$
\cite{nist-quadrature}. Bei einem Simplexpolynom vom Gesamtgrad $d$ werden die
Grade nach der Tetraederabbildung höchstens $d+2,d+1,d$.
Vier Punkte je Richtung geben 64 positive Punkte und garantierten Gesamtgrad
fünf; auf Flächen ergeben drei Punkte je Richtung neun Punkte und Grad vier.
Die physischen Gewichte summieren sich auf $V$ bzw. $A$ und tragen m$^3$ bzw.
m$^2$. Das ist ein verifizierbarer Referenzalgorithmus, keine minimale Regel.
Produkte $N_aN_b$ sind Grad vier, Produkte zweier Gradienten Grad zwei.
Nichtpolynomiale Lasten verlangen eine eigene Quadraturstudie.

\section{Randtopologie und Lastmomente}
Jede Außenfläche muss genau einem Volumenelement gehören. Die Normale wird
gegen die gegenüberliegende Ecke nach außen orientiert. Eine umgekehrte
importierte Dreiecksreihenfolge verändert deshalb keinen vorgegebenen äußeren
Fluss. Innere, fremde und nichtmannigfaltige Flächen werden diagnostiziert;
Teilrandnetze sind für Solves zulässig, der Wärmeindikator verlangt vollständige
Randabdeckung. Konforme Materialinterfaces besitzen zwei äußere Normalen
$\mathbf n_+=-\mathbf n_-$ und keine zusätzlichen Interface-DOFs.

Eine konstante Triangle3-Last hat Gewicht $A/3$ je Ecke. Bei Triangle6 gilt
$\int_F N_i=0$ an Ecken und $A/3$ an Kantenmitten. Dies folgt direkt aus
Gleichung~\ref{eq:simplex-moment} in zwei Dimensionen. Null Ecklast bedeutet
nicht, dass ein Eckknoten aus einer Elektrode entfernt werden darf.
Volumenlasten auf Tet10 haben $-V/20$ je Ecke und $V/5$ je Kantenknoten;
negative konsistente Einträge verletzen nicht die positive Gesamtlast.

\section{Ausgearbeitetes Einheitsbeispiel}\label{sec:unit-tet}
Für Ecken $(0,0,0),(1,0,0),(0,1,0),(0,0,1)$ in m ist $V=1/6$ m$^3$ und
\begin{equation}\label{eq:unit-G}
\mathbf G=\begin{pmatrix}-1&1&0&0\\-1&0&1&0\\-1&0&0&1\end{pmatrix}\ {\rm m}^{-1}.
\end{equation}
$T=x$ (numerische SI-Koordinaten) wird durch Tet4 exakt reproduziert.
Für $k=2$ W/(m K) ist $\mathbf q=(-2,0,0)$ W/m$^2$ und
$\mathbf T_e^T\mathbf K_T\mathbf T_e=1/3$ W K.
Tet10 reproduziert $T=x^2$; dann
$\int_K\nabla T^T\mathbf k\nabla T=8\int_K x^2=8/60=2/15$ W K.
Die letzte Rechnung bezieht sich auf numerisch in Metern eingesetzte
Koordinaten; die Polynomkoeffizienten haben die nötigen K/m bzw. K/m$^2$.
Sie prüft quadratische Feldenergie statt nur einen Schwerpunktwert.

\subsection*{Implementierung und Nachweis.}
{\small\raggedright
\path{src/fem_toolbox/elements/tet4.py},
\path{src/fem_toolbox/elements/tet10.py},
\path{src/fem_toolbox/elements/quadrature.py};
\path{tests/unit/test_tet4_geometry.py},
\path{tests/unit/test_tet10_geometry.py},
\path{tests/unit/test_tet10_operators.py},
\path{tests/convergence/test_tet10_interpolation.py}. Die unabhängige
Monomintegration und alle 24 Eckpermutationen prüfen die Abbildung.

\par}
